\documentclass[11pt]{letter}
\usepackage[margin=2.5cm]{geometry}
\usepackage[T1]{fontenc}
\usepackage{hyperref}
\signature{Author Name\\Affiliation\\Email}
\address{Author Name\\Affiliation\\City, Country\\Email}
\begin{document}
\begin{letter}{The Editors, \emph{Explorations in Economic History}}
\opening{Dear Editors,}

I am pleased to submit the manuscript entitled ``Network Exclusion and State Collapse: From Maritime Isolation to Technological Access Denial in the Long Run of History'' for consideration by \emph{Explorations in Economic History}.

Using an AI-assisted exploratory dataset of 96 historical polities rather than a public-data measurement dataset, we distinguish between deliberate closure (policy-based trade bans, \emph{sakoku}, bloc membership) and what we term \emph{technical network exclusion}: involuntary disconnection from the dominant exchange network of an era due to geographic or technological constraints. The manuscript presents four features that we believe are of interest to the journal's readership.

First, a post hoc sensitivity analysis reports that reclassifying 7 technically excluded polities changes the one-sided Fisher exact p-value for closure and conquest from 0.187 to 0.020 within the current codings, while the core stock--flow odds ratio remains stable (OR $= 1.774$).

Second, we analyze conditional closure---cases where polities maintained selective technology channels while restricting broader trade. Examples include Tokugawa Japan's \emph{rangaku} through Dejima, Qing China's Canton system, Joseon Korea's tributary trade, and the Soviet Union's bloc-internal technology sharing. The mixed outcomes of these cases suggest that the open/closed dichotomy is insufficient; the conditions under which selective channels mitigate technology gaps merit further investigation.

Third, the closure-type comparison shows a 100.0\% conquest rate among the technically excluded candidates, but the category rates are not strictly monotonic. We therefore present technology-flow disruption as a hypothesis rather than an identified mechanism.

Fourth, the hypothesis may generalize beyond geographic isolation: as the dominant network shifts from sea lanes to semiconductors, AI, and advanced robotics, states structurally excluded from these platforms may face analogous vulnerabilities.

The manuscript includes the analysis tables, figures, and a complete Supplementary Table~S1 dataset. [Before submission, please confirm that this work has not been published or submitted elsewhere.]

We suggest the following reviewers based on their expertise in quantitative economic history:
\begin{itemize}
\item Prof.\ J\"org Baten (University of T\"ubingen) --- cliometric methods and long-run development
\item Prof.\ Stephen Broadberry (University of Oxford) --- comparative economic history and GDP estimation
\item Prof.\ Nathan Nunn (University of British Columbia) --- long-run persistence and historical institutions
\end{itemize}

\closing{Sincerely,}
\end{letter}
\end{document}